\section{Ergebnisse}
\label{sec:results}

\subsection{Überblick über die Evaluationsergebnisse}
\label{sec:results_overview}

Die in Kapitel~\ref{sec:validation} beschriebene Evaluation liefert
unterschiedliche Evidenz für die technische Realisierung und die
Tragfähigkeit des entwickelten Architekturkonzepts. Die Ergebnisse werden
daher nicht als einzelner Erfolgswert zusammengefasst, sondern entlang der
untersuchten Verifikations- und Validierungsebenen dargestellt.

Die automatisierte Softwareverifikation bestätigte zunächst die im Prototyp
implementierten deterministischen Contracts und Zustandsübergänge. Neben
fokussierten Tests einzelner Komponenten und Architekturgrenzen wurden
wiederholt vollständige Regressionen durchgeführt. Damit konnte insbesondere
geprüft werden, dass Änderungen an einzelnen Verarbeitungsschritten bereits
implementierte Upstream- und Downstream-Verträge nicht unbeabsichtigt
verletzten. Die konkreten Ergebnisse dieser Prüfungen werden in
Abschnitt~\ref{sec:results_software} zusammengefasst.

Die formative End-to-End-Erprobung zeigte darüber hinaus, dass technisch
korrekte Einzelkomponenten allein nicht ausreichen, um einen kontrollierten
Engineering Information Flow sicherzustellen. Während der verbundenen
Ausführung wurden mehrere Grenzen zwischen Evidence, semantischer
Interpretation, Human Authority und Modellbildung präzisiert. Besonders
relevant waren dabei Beobachtungen, aus denen Änderungen oder
Konkretisierungen des Architekturkonzepts hervorgingen. Reine
Implementierungsfehler wurden ebenfalls dokumentiert und durch Korrektur und
Retest behandelt, stehen jedoch nicht im Mittelpunkt der Ergebnisdarstellung.

Die für das Architekturkonzept relevanten Findings betreffen insbesondere die
Trennung von source-grounded Evidence und Persona Interpretation, die
explizite Autorisierung semantischer Relationships, die Unterscheidung von
Engineering Meaning, Model Placement und Target-Model Formulation sowie die
Begrenzung automatisierter beziehungsweise LLM-gestützter Verarbeitung an
expliziten Authority-Grenzen. Diese Ergebnisse werden in
Abschnitt~\ref{sec:results_formative} dargestellt. Eine vollständige
Aufbereitung der im Entwicklungs- und Evaluationsprozess dokumentierten
Findings mit Beobachtung, Reaktion und Status wird ergänzend im Anhang
bereitgestellt. Dort werden diejenigen Findings gesondert hervorgehoben, die
zu einer Änderung oder Präzisierung des Architekturkonzepts geführt haben.

Neben dem formativen Einzelpfad wurde die Multi-Source-Verarbeitung separat
untersucht. Nach Korrektur einer zunächst blockierenden
Identitäts- und Provenance-Problematik konnte der kontrollierte
Multi-Source-End-to-End-Pfad vollständig durchlaufen werden. Damit lag für
den untersuchten Prototyp ein erfolgreicher Verarbeitungspfad vor, bei dem
mehrere getrennte Engineering Sources source-lokal verarbeitet und erst nach
ihrer jeweiligen Autorisierung in die projektweite Modellbildung überführt
wurden. Die entsprechenden Ergebnisse werden in
Abschnitt~\ref{sec:results_multisource} zusammengefasst.

Den abschließenden Nachweis des konsolidierten Verarbeitungskonzepts bildet
der Referenzlauf des Projekts~930444. Dabei wurde der vollständige
Engineering Information Flow von den registrierten Engineering- und Context
Sources bis zum freigegebenen SysML-v2-Artefakt ausgeführt. Entlang dieses
Workflows wurden die relevanten Evidence-, Decision-, Provenance-,
Authority-, Modell-, Validierungs- und Release-Zustände persistiert.

Der dabei entstandene Evidenzbestand ist wesentlich umfangreicher als im
Haupttext sinnvoll wiedergegeben werden kann. Die Traceability wird deshalb
in Abschnitt~\ref{sec:results_traceability} exemplarisch anhand eines
einzelnen fachlichen Informationspfads dargestellt. Ausgewählt wird ein Fall,
bei dem ein maschinell vorbereiteter Review-Zustand durch eine menschliche
Entscheidung verändert und anschließend über die Modellbildung bis in das
generierte SysML-v2-Artefakt weitergeführt wurde. Dadurch lässt sich nicht
nur die Source-Provenance, sondern auch die Wirkung der Human Authority und
die Weiterverarbeitung des tatsächlich autorisierten Informationszustands
nachvollziehen.

Die vollständige Traceability des Referenzprojekts umfasst demgegenüber alle
zugelassenen Informationspfade und die zugehörigen Zwischen- und
Entscheidungszustände. Eine vollständige Wiedergabe würde den Umfang der
Arbeit deutlich erhöhen, ohne für jeden einzelnen Pfad einen zusätzlichen
Erkenntnisgewinn zu liefern. Die Ergebnisdarstellung konzentriert sich daher
auf die Gesamtstruktur der erzeugten Evidenz und auf einen vollständig
nachvollziehbaren exemplarischen Durchgang.

Schließlich wurde das erzeugte SysML-v2-Artefakt technisch mit einer externen
SysML-v2-Werkzeugumgebung validiert und anschließend einer getrennten Human
Review und Release Decision unterzogen. Die technische Validierung und die
fachliche Freigabe blieben damit auch im ausgeführten End-to-End-Workflow
getrennte Zustände. Die entsprechenden Ergebnisse werden in
Abschnitt~\ref{sec:results_validation} dargestellt.

Die formative Verifikation erzeugte neben den im Haupttext dargestellten
Ergebnissen eine größere Zahl technischer, semantischer sowie
bedienungsbezogener Findings. Die Auswahl für die Ergebnisdarstellung erfolgt
nicht anhand der internen Schweregrad- oder Finding-Kategorie, sondern anhand
ihrer Bedeutung für den Untersuchungsgegenstand. Befunde werden im Haupttext
einzeln dargestellt, wenn sie eine angenommene Verantwortungs-, Informations-,
Provenance- oder Authority-Grenze des Architekturkonzepts infrage stellten oder
eine wesentliche Ergebnisgrenze sichtbar machten.

Reine Implementierungsfehler wurden, soweit sie den weiteren
Verifikationspfad blockierten, korrigiert und am betroffenen Verarbeitungsschritt
erneut geprüft, ohne daraus einen eigenständigen Architekturbeitrag abzuleiten.
Nach Abschluss des validierten Proof-of-Concept-Umfangs verbleibende
nicht blockierende semantische, integrative und bedienungsbezogene Findings
wurden nicht weiter optimiert. Sie bleiben als dokumentierte Limitationen und
Ausgangspunkte für eine Weiterentwicklung erhalten.

Die vollständige Finding-Historie einschließlich Erstbeobachtung, Reaktion,
Retest und finalem Status ist in
Anhang~\ref{app:findings_register} dokumentiert. Historische Fehlerzustände
werden dort nicht durch spätere erfolgreiche Retests ersetzt.

Tabelle~\ref{tab:results_overview} fasst die wesentlichen Ergebnisbereiche
zusammen.

\begin{table}[htbp]
    \centering
    \small
    \caption{Überblick über die Evaluationsergebnisse}
    \label{tab:results_overview}
    \begin{tabular}{p{3.8cm}p{5.0cm}p{4.7cm}}
        \hline
        \textbf{Evaluationsbereich} &
        \textbf{Zentraler Prüfgegenstand} &
        \textbf{Ergebnischarakter} \\
        \hline

        Software- und Contract-Verifikation &
        deterministische Implementierungs- und
        Authority-Verträge &
        technische Verifikation der implementierten
        Architekturmechanismen \\

        Formative End-to-End-Erprobung &
        Zusammenspiel von Evidence, Interpretation,
        Human Authority und Modellbildung &
        Identifikation und Präzisierung relevanter
        Architekturgrenzen \\

        Multi-Source-Verifikation &
        source-lokale Verarbeitung, Provenance und
        projektweite Modellbildung &
        erfolgreicher kontrollierter
        Multi-Source-End-to-End-Pfad \\

        Referenzlauf 930444 &
        vollständiger Informations- und
        Authority-Pfad des konsolidierten Prototyps &
        persistierter Evidence-, Decision-,
        Provenance- und Traceability-Artefaktsatz \\

        SysML-v2-Validierung und Release &
        technische Verarbeitbarkeit und getrennte
        fachliche Freigabe &
        externe technische Validierung und explizite
        Human Release \\
        \hline
    \end{tabular}
\end{table}

Die folgenden Abschnitte stellen diese Ergebnisse im Einzelnen dar. Die
wissenschaftliche Interpretation der Befunde und ihre Rückbindung an die
Forschungsfragen erfolgen davon getrennt in Kapitel~\ref{sec:discussion}.

\subsection{Ergebnisse der Software- und Contract-Verifikation}
\label{sec:results_software}

Die automatisierte Verifikation begleitete die prototypische Realisierung über
den gesamten Entwicklungszeitraum. Nach Änderungen an Datenstrukturen,
Zustandsübergängen oder Architekturgrenzen wurden zunächst fokussierte
Unit-, Contract- und Integrationstests ausgeführt. Anschließend wurden die
betroffenen Workflowpfade sowie die vollständige Repository-Regression erneut
geprüft.

Dabei wurden insbesondere die für das Architekturkonzept relevanten
deterministischen Eigenschaften abgedeckt. Hierzu gehörten eindeutige
Identitäten und Fingerprints, immutable Zustände und Revisionen,
Provenance-Bindungen, zulässige Authority-Übergänge, die Persistenz von
Entscheidungen sowie fail-closed Verhalten bei fehlenden oder nicht
autorisierten Voraussetzungen.

Die fokussierten Tests erwiesen sich insbesondere während der formativen
End-to-End-Erprobung als relevant. Wurde an einer verbundenen
Verarbeitungsgrenze eine Abweichung festgestellt, konnte die zugehörige
Korrektur zunächst isoliert gegen den betroffenen Contract geprüft werden,
bevor derselbe Verarbeitungsschritt erneut im End-to-End-Workflow ausgeführt
wurde. Zusätzlich wurde durch die vollständige Regression kontrolliert, dass
die Korrektur bereits etablierte Verarbeitungspfade nicht unbeabsichtigt
veränderte.

Für einen vor Abschluss der Multi-Source-Erweiterung akzeptierten
End-to-End-Baseline-Stand umfasste die vollständige Repository-Regression
6100 erfolgreich ausgeführte Tests. Auch die nachfolgenden
Multi-Source-Korrekturen wurden durch fokussierte Tests der betroffenen
Identitäts-, Provenance-, Handoff- und Modellierungsverträge sowie durch
kombinierte Regressionen abgesichert. Die jeweilige technische
Repository-Konsistenzprüfung blieb dabei ohne Befund.

Tabelle~\ref{tab:results_software_verification} fasst die beobachteten
Verifikationsergebnisse auf den drei Prüfebenen zusammen.

\begin{table}[htbp]
    \centering
    \small
    \caption{Ergebnisse der automatisierten Softwareverifikation}
    \label{tab:results_software_verification}
    \begin{tabular}{p{3.6cm}p{6.0cm}p{3.8cm}}
        \hline
        \textbf{Prüfebene} &
        \textbf{Beobachtetes Ergebnis} &
        \textbf{Status} \\
        \hline

        Unit / Contract &
        deterministische Daten-, Identitäts-, Persistenz- und
        Authority-Verträge konnten fokussiert geprüft werden &
        bestanden \\

        Integration / Boundary &
        Übergänge zwischen Processing, Human Authority,
        Modellbildung, Generierung und Release wurden in
        fokussierten Regressionen geprüft &
        bestanden nach ggf. erforderlicher Korrektur und Retest \\

        Vollständige Regression &
        bereits akzeptierte Verträge wurden nach Änderungen
        erneut gemeinsam ausgeführt; für den validierten
        End-to-End-Baseline-Stand wurden 6100 Tests erfolgreich
        ausgeführt &
        bestanden \\
        \hline
    \end{tabular}
\end{table}

Während der verbundenen Erprobung traten trotz bestandener isolierter Tests
mehrere Integrations- und Implementierungsfehler auf. Diese Befunde wurden
nicht durch die Regression verdeckt, sondern als eigene Evaluationsevidenz
behandelt, korrigiert und am jeweils betroffenen Verarbeitungsschritt erneut
geprüft. Nur ein Teil dieser Befunde offenbarte ein konzeptionelles
Architekturproblem; die übrigen betrafen die konkrete prototypische
Implementierung.

Die Testanzahl wird daher nicht als Maß für die Qualität oder wissenschaftliche
Validität des Architekturkonzepts interpretiert. Das Ergebnis der
automatisierten Verifikation besteht vielmehr darin, dass die spezifizierten
deterministischen Contracts des untersuchten Prototyps im jeweils geprüften
Softwarestand ausführbar und regressionsgesichert waren. Die darüber
hinausgehenden Erkenntnisse aus dem verbundenen Engineering-Workflow werden
im folgenden Abschnitt dargestellt.

\subsection{Ergebnisse der formativen End-to-End-Erprobung}
\label{sec:results_formative}

Die formative End-to-End-Erprobung führte sowohl zu technischen
Implementierungsbefunden als auch zu Findings mit unmittelbarer Bedeutung für
das Architekturkonzept. Im Folgenden werden ausschließlich diejenigen Befunde
einzeln dargestellt, durch die eine zunächst angenommene Verantwortungs-,
Informations-, Provenance- oder Authority-Grenze angepasst oder präzisiert
werden musste. Die vollständige Finding-Historie ist in
Anhang~\ref{app:findings_register} dokumentiert.

\paragraph{Trennung von Engineering Evidence und Verarbeitungskontext}

In frühen verbundenen Verarbeitungsläufen wurden Teile des für die
Orchestrierung, Aufgabenbeschreibung oder Interpretation bereitgestellten
Kontexts teilweise als positive Engineering Information behandelt. Damit
konnten Informationen in den fachlichen Verarbeitungspfad gelangen, obwohl
sie nicht durch die eigentliche Engineering Source belegt waren.

Die Verarbeitung wurde daraufhin so abgegrenzt, dass Engineering Evidence
eine explizite source-grounded Grundlage benötigt. Kontextinformationen
dürfen weiterhin zur Interpretation, Terminologieauflösung oder
Klassifikation beitragen, erhalten dadurch jedoch keine eigenständige
Evidence-Authority. Diese Trennung wurde anschließend als durchgängige
Informationsgrenze in die Architektur übernommen.

\paragraph{Gemeinsame Subjects vor der Persona-Interpretation}

Ein weiterer Befund betraf die Multi-Perspektiven-Verarbeitung. Werden
Personas bereits zur Bildung fachlicher Subjects eingesetzt, können
unterschiedliche Perspektiven voneinander abweichende Subject-Populationen
erzeugen. Unterschiede zwischen den Persona-Ergebnissen lassen sich dann
nicht mehr eindeutig als unterschiedliche Interpretationen desselben
Engineering-Gegenstands bewerten.

Die Verarbeitung wurde deshalb in zwei Verantwortungen getrennt.
Source-grounded Evidence wird zunächst in gemeinsame Canonical Engineering
Subjects überführt. Erst anschließend interpretieren mehrere Personas
denselben fachlichen Gegenstand. Consensus, gemeinsame Unsicherheit und
tatsächliche Interpretationsvarianz können dadurch auf einen gemeinsamen
Referenzgegenstand zurückgeführt und dem Human Review getrennt bereitgestellt
werden.

\paragraph{Umfang autorisierter Engineering Information}

Während der verbundenen Weiterverarbeitung zeigte sich zudem, dass eine
Autorisierung einzelner fachlicher Subjects nicht ausreicht, wenn im Human
Review ebenfalls semantische Beziehungen zwischen diesen Subjects
entschieden werden. In einem Zwischenstand wurden akzeptierte Relationships
zwar im Review geführt, waren jedoch nicht Bestandteil des autoritativen
Handoffs zur nachfolgenden Modellableitung.

Der Authority-Zustand wurde daher erweitert. Die nachgelagerte Modellbildung
erhält neben den freigegebenen fachlichen Subjects auch die explizit
akzeptierten semantischen Relationships. Damit wird nicht erst während der
Modellbildung erneut entschieden, welche zuvor geprüften fachlichen
Beziehungen gültig sind.

\paragraph{Engineering Meaning und Zielmodellbildung}

Der deutlichste strukturelle Befund trat beim Übergang von autorisierter
Engineering Information zur Modellbildung auf. Eine direkte Ableitung vom
fachlich freigegebenen Inhalt zur konkreten SysML-v2-Repräsentation erwies
sich als zu grob. Die fachliche Bedeutung einer Information legt nicht
eindeutig fest, an welcher Stelle des Zielmodells sie repräsentiert oder wie
sie für die konkrete Zielsprache formuliert werden soll.

Die Verarbeitung wurde deshalb in getrennte Entscheidungsstufen
aufgeteilt:

\begin{center}
Engineering Meaning\\
$\downarrow$\\
Model Placement / Representation\\
$\downarrow$\\
Target-Model Formulation\\
$\downarrow$\\
SysML-v2-Serialisierung.
\end{center}

Model Placement entscheidet damit über die strukturelle Einordnung bereits
autorisierter Engineering Information. Die nachfolgende
Target-Model-Formulation behandelt die modellgerechte Ausprägung innerhalb
dieser autorisierten Einordnung. Beide Entscheidungen bleiben von der
ursprünglichen fachlichen Autorisierung unterscheidbar.

\paragraph{Deterministische Model Assembly}

Ein weiterer verbundener Testlauf zeigte, dass die Model Assembly in einem
Zwischenstand erneut LLM-basierte Model-Placement-Perspektiven verwendete, um
die Repräsentation bereits akzeptierter Relationships zu bestimmen. Dadurch
entstand innerhalb einer eigentlich nachgelagerten Modellbildungsstufe eine
neue semantische Entscheidungsinstanz.

Die Verantwortung der Model Assembly wurde daraufhin eingeschränkt.
Eindeutig autorisierte und durch die Modellierungsregeln abbildbare
Information wird deterministisch übernommen. Kann eine fachlich akzeptierte
Relationship nicht eindeutig durch die zugelassenen Modellierungsregeln
repräsentiert werden, wird sie nicht durch einen impliziten AI-Fallback
aufgelöst, sondern als explizit zu entscheidender Zustand an die zuständige
Human Authority weitergegeben.

Tabelle~\ref{tab:results_architecture_findings} fasst die für die
Architekturentwicklung wesentlichen Befunde zusammen.

\begin{table}[htbp]
    \centering
    \small
    \caption{Architekturrelevante Befunde der formativen End-to-End-Erprobung}
    \label{tab:results_architecture_findings}
    \begin{tabular}{p{4.0cm}p{4.8cm}p{4.8cm}}
        \hline
        \textbf{Beobachtung} &
        \textbf{Architektonische Reaktion} &
        \textbf{Resultierender Zustand} \\
        \hline

        Verarbeitungskontext konnte als fachliche Information erscheinen &
        Trennung von Evidence Basis und Context Basis &
        nur source-grounded Information kann Engineering Evidence begründen \\

        Persona-Verarbeitung erzeugte voneinander abweichende
        Subject-Populationen &
        gemeinsame Subject-Bildung vor der Persona-Interpretation &
        Personas interpretieren identische Canonical Engineering Subjects \\

        akzeptierte Relationships fehlten im autoritativen
        Modellierungs-Handoff &
        Erweiterung der autorisierten Engineering Information &
        Subjects und akzeptierte Relationships werden gemeinsam weitergegeben \\

        Engineering Meaning wurde zu direkt auf die
        Zielrepräsentation abgebildet &
        Trennung von Meaning, Placement und Formulation &
        getrennte fachliche und modellbezogene Authority-Stufen \\

        Model Assembly traf erneut semantische Entscheidungen &
        deterministische Assembly mit expliziter Eskalation &
        keine implizite semantische Neuentscheidung während der Assembly \\
        \hline
    \end{tabular}
\end{table}

Neben diesen architekturrelevanten Befunden wurden weitere technische,
integrative und bedienungsbezogene Findings dokumentiert. Technische Defekte
wurden korrigiert und durch fokussierte Tests sowie Retests abgesichert,
führten jedoch nicht notwendigerweise zu einer Änderung des
Architekturkonzepts. Beobachtungen zur Bedienbarkeit betreffen vor allem die
prototypische Softwareausprägung und werden in den Limitationen und im
Ausblick erneut aufgegriffen.

Eine vollständige Aufbereitung der während der Evaluation dokumentierten
Findings enthält der Anhang. Dort werden für jedes Finding die beobachtete
Abweichung und die darauf erfolgte Reaktion beziehungsweise Korrektur
angegeben. Die in diesem Abschnitt dargestellten Findings werden dort als
architekturrelevant gekennzeichnet.

\subsection{Ergebnisse der Multi-Source-Verifikation}
\label{sec:results_multisource}

Die Multi-Source-Verifikation untersuchte, ob mehrere getrennte Engineering
Sources innerhalb eines gemeinsamen Projekts verarbeitet werden können, ohne
ihre jeweilige Herkunft und fachliche Authority zu verlieren. Hierfür wurden
mehrere synthetische Legacy-Dokumente separat registriert und zunächst
source-lokal durch Processing und Human Engineering Review geführt.

Der erste verbundene Multi-Source-Durchlauf konnte nicht vollständig
abgeschlossen werden. An der projektweiten Weiterverarbeitung trat eine
Identitätskollision zwischen Artefakten unterschiedlicher Processing Runs
auf. Dadurch war nicht mehr in jedem Fall eindeutig bestimmbar, welchem
Source- und Processing-Kontext ein nachgelagertes Artefakt zugeordnet werden
musste. Der betroffene Verarbeitungspfad wurde entsprechend der
fail-closed Logik nicht weitergeführt. Die vollständige Historie dieses Befunds mit Erstbeobachtung,
Architekturkorrektur und erfolgreichem Retest ist ebenfalls in
Anhang~\ref{app:findings_register} dokumentiert.

Die Korrektur bestand nicht in einer Zusammenführung der source-lokalen
Zustände. Stattdessen wurde die projektweite Identität von Artefakten um
ihren Processing-Kontext erweitert. Zusätzlich wurden die Übergaben zwischen
source-lokaler Engineering Authority und projektweiter Modellbildung an die
jeweils konkreten Source-, Run-, Review- und Fingerprint-Zustände gebunden.
Damit blieb die Herkunft auch nach Eintritt in die projektweite Verarbeitung
explizit erhalten.

Nach dieser Korrektur wurde der kontrollierte Multi-Source-Pfad erneut
ausgeführt. Der Retest konnte vollständig durchlaufen werden. Dabei wurden
die Engineering Sources zunächst unabhängig voneinander verarbeitet und
Human-reviewed. Die daraus resultierende autorisierte Engineering
Information blieb source-lokal erhalten und wurde anschließend über einen
expliziten Project-Fit-Schritt für die gemeinsame Modellbildung zugelassen.

Der erfolgreiche Verarbeitungspfad entsprach damit folgender Struktur:

\begin{center}
mehrere Engineering Sources\\
$\downarrow$\\
source-lokales Processing\\
$\downarrow$\\
source-lokaler Human Engineering Review\\
$\downarrow$\\
source-lokale Approved Engineering Information\\
$\downarrow$\\
Project Fit / Admission\\
$\downarrow$\\
gemeinsamer Model Candidate Set\\
$\downarrow$\\
Model Authority\\
$\downarrow$\\
gemeinsames Systemmodell.
\end{center}

Ein wesentliches Ergebnis ist, dass die projektweite Verarbeitung keine
synthetisch zusammengeführte Approved Engineering Information erzeugt.
Stattdessen werden die einzelnen autorisierten Source-Zustände getrennt
weitergeführt und gemeinsam als zugelassene Eingabe für die Modellbildung
verwendet. Dadurch bleibt nachvollziehbar, welche Source einen bestimmten
fachlichen Beitrag geliefert hat, auch wenn mehrere Sources anschließend in
dasselbe Projektmodell eingehen.

Gleichzeitig entsteht für den zugelassenen Multi-Source-Zustand genau ein
gemeinsamer Model Candidate Set. Die Source-Grenzen bestimmen damit die
Provenance der zugrunde liegenden Engineering Information, erzwingen jedoch
keine getrennten Zielmodelle. Mehrere Sources können zum selben
Modellierungsgegenstand beitragen, ohne dass ihre jeweilige Herkunft dabei
verloren geht.

Tabelle~\ref{tab:results_multisource} fasst die zentralen Ergebnisse der
Multi-Source-Verifikation zusammen.

\begin{table}[htbp]
    \centering
    \small
    \caption{Ergebnisse der Multi-Source-Verifikation}
    \label{tab:results_multisource}
    \begin{tabular}{p{4.4cm}p{8.6cm}}
        \hline
        \textbf{Prüfgegenstand} &
        \textbf{Beobachtetes Ergebnis} \\
        \hline

        Source-Identität &
        getrennt registrierte Engineering Sources blieben über die
        source-lokale Verarbeitung eindeutig unterscheidbar \\

        Processing-Provenance &
        nach Korrektur der initialen Identitätskollision konnten
        projektweite Referenzen eindeutig auf den jeweiligen
        Processing-Kontext zurückgeführt werden \\

        Human Authority &
        nur zuvor source-lokal geprüfte und autorisierte Engineering
        Information wurde für die projektweite Weiterverarbeitung zugelassen \\

        Projektweite Zulassung &
        die Aufnahme in die gemeinsame Modellbildung erfolgte über einen
        expliziten Project-Fit-Zustand mit Bindung an die zugrunde liegende
        Source- und Authority-Evidence \\

        Source-lokale Provenance &
        die projektweite Verarbeitung erzeugte keine synthetisch
        zusammengeführte Approved Engineering Information; die einzelnen
        autorisierten Source-Zustände blieben erhalten \\

        Gemeinsame Modellbildung &
        die zugelassenen source-lokalen Informationen konnten in einen
        gemeinsamen Model Candidate Set und den nachfolgenden
        Modellbildungsprozess überführt werden \\

        End-to-End-Ergebnis &
        der korrigierte Multi-Source-Verarbeitungspfad wurde im Retest
        vollständig durchlaufen \\
        \hline
    \end{tabular}
\end{table}

Die Multi-Source-Verifikation zeigte damit im untersuchten Prototyp, dass
source-lokale Provenance und gemeinsame projektweite Modellbildung
miteinander vereinbar sind. Die Verifikation bezieht sich auf den
kontrollierten Testfall und die implementierten Verarbeitungspfade. Sie
belegt keine allgemeine Vollständigkeit oder Konfliktfreiheit beliebiger
industrieller Multi-Source-Datensätze.

Der abschließende Referenzlauf in
Abschnitt~\ref{sec:results_traceability} erweitert diese Betrachtung um die
persistierte End-to-End-Traceability eines vollständig durchgeführten
Projektlaufs.

\subsection{Referenzlauf 930444 und exemplarische End-to-End-Traceability}
\label{sec:results_traceability}

Den abschließenden End-to-End-Nachweis des konsolidierten Prototyps bildet
der Referenzlauf des Projekts~930444. Der Lauf wurde als kontrollierter
Testlauf über die reguläre Benutzeroberfläche durchgeführt. Dabei wurden die
Verarbeitungsschritte nicht lediglich isoliert aufgerufen, sondern entlang
des implementierten Workflows bis zur technischen Validierung und
anschließenden Human Release Decision durchlaufen.

Für den Referenzlauf waren insgesamt sechs Sources registriert. Vier davon
wurden als \texttt{engineering\_source} verarbeitet, zwei weitere als
\texttt{context\_only}. Die vier Engineering Sources enthielten jeweils
unterschiedliche Ausschnitte desselben beispielhaften
Remote-Microscopy-Kontexts: Stakeholder, User Workflow, System Requirements
und Architecture Notes. Die Context Sources beschrieben den übergeordneten
Produktkontext. Die vollständigen verwendeten Testdateien und ihre
Zuordnung zu den persistierten Source-Identitäten sind in
Anhang~\ref{app:traceability_930444_inputs} wiedergegeben.

Im vollständigen Lauf wurden elf source-grounded Evidence Items und zwanzig
Canonical Engineering Subjects persistiert. Die Engineering Sources wurden
source-lokal verarbeitet und jeweils einem Human Engineering Review
zugeführt. Aus den finalisierten Reviews entstanden insgesamt zwanzig
Approved Inputs und vier Project-Fit-Zustände. Für die Modellbildung wurden
zwanzig Human Model Placement Decisions persistiert. Der abschließende
Modellzustand umfasste zwanzig Elemente und sechs Relationships. Für den
generierten Output wurden 26 Traceability Entries erzeugt.

\paragraph{LLM-Nutzung des Referenzlaufs}

Neben den fachlichen und technischen Ergebnisartefakten wurde für den
Referenzlauf auch die Nutzung der angebundenen Sprachmodelle ausgewertet.
Hierzu wurden die in den persistierten LLM-Responses enthaltenen
Usage-Informationen anhand ihrer eindeutigen Response-Identitäten
zusammengeführt. Mehrfach innerhalb des Evidence-Satzes persistierte Kopien
derselben Response wurden dabei nur einmal berücksichtigt. Für den
Referenzlauf lassen sich damit 35 unterschiedliche LLM-Aufrufe mit insgesamt
69\,507 Input- und 21\,693 Output-Tokens rekonstruieren. Tabelle~
\ref{tab:results_930444_token_usage} zeigt die Verteilung auf die verwendeten
Modelle.

\begin{table}[htbp]
    \centering
    \small
    \caption{LLM-Nutzung im Referenzlauf 930444}
    \label{tab:results_930444_token_usage}
    \begin{tabular}{lrrrr}
        \hline
        \textbf{Modell} &
        \textbf{Requests} &
        \textbf{Input-Tokens} &
        \textbf{Output-Tokens} &
        \textbf{Gesamt} \\
        \hline
        \texttt{gpt-5.4-mini} & 33 & 63\,392 & 17\,333 & 80\,725 \\
        \texttt{gpt-5.5}      & 2  & 6\,115  & 4\,360  & 10\,475 \\
        \hline
        \textbf{Gesamt}       & \textbf{35} & \textbf{69\,507} &
        \textbf{21\,693} & \textbf{91\,200} \\
        \hline
    \end{tabular}
\end{table}

Die Summe von 91\,200 Tokens beschreibt ausschließlich die im archivierten
Referenzlauf eindeutig identifizierbaren LLM-Responses. Der zusätzlich
herangezogene, tagesaggregierte Usage-Export des API-Providers weist für den
11.~September~2026 weitere API-Aktivität desselben Projekts aus und ist daher
nicht vollständig dem Referenzlauf zurechenbar. Er wird lediglich zur
Plausibilisierung der persistierten Usage-Daten verwendet. Eine Bewertung der
Effizienz erfolgt in Abschnitt~\ref{sec:discussion_limitations}.

Der dabei entstandene Evidence-, Decision-, Provenance- und
Traceability-Artefaktsatz ist wesentlich umfangreicher, als im Haupttext
vollständig wiedergegeben werden kann. Die End-to-End-Traceability wird
deshalb im Folgenden exemplarisch anhand ausgewählter Originalartefakte
dargestellt. Als durchgehender fachlicher Gegenstand wurde
\emph{Microscope Workstation} ausgewählt. An diesem Beispiel lassen sich
Source Evidence, Canonical Subject, unterschiedliche
Persona-Interpretationen, Human Engineering Review, Approved Engineering
Information, Model Placement, Internal Engineering Model, Model Quality
Decision und die spätere Zuordnung zum generierten SysML-v2-Artefakt
gemeinsam verfolgen.

Die Auswahl dieses Gegenstands dient ausschließlich dazu, die persistierte
Verarbeitungskette an einem zusammenhängenden Beispiel sichtbar zu machen.
Einzelne weitere Artefaktarten, die für \emph{Microscope Workstation} nicht
als eigener Entscheidungsschritt vorliegen, werden ergänzend anhand anderer
Gegenstände des gleichen Referenzlaufs gezeigt. Die für die Rekonstruktion
verwendeten Originalauszüge sind in
Anhang~\ref{app:traceability_930444_workstation} und
Anhang~\ref{app:traceability_930444_additional} zusammengeführt.


\paragraph{Fehlgeschlagener Processing Attempt und Retry}

Die Verarbeitung von \texttt{SRC-000005} erfolgte innerhalb von
\texttt{RUN-000001}. Der erste Processing Attempt
\texttt{ATT-000001} wurde während der agentischen
Informationsverarbeitung beendet und als fehlgeschlagener Zustand
persistiert. Abbildung~\ref{fig:result_930444_failed_attempt} zeigt den
entsprechenden Ausschnitt aus der Event-Historie.

\begin{figure}[htbp]
    \centering
    \begin{minipage}{0.94\textwidth}
        \begin{evidencebox}
event_id:         EVT-000003
previous_state:   running
next_state:       failed
processing_stage: agentic_ingestion
event_type:       run_failed
attempt_id:       ATT-000001
reason_code:      shared_evidence_interpretation_failed

EVT-000004:
previous_state:   failed
next_state:       running
event_type:       retry_started
attempt_id:       ATT-000002
reason_code:      agentic_ingestion_retry_started
        \end{evidencebox}
    \end{minipage}
    \caption{Persistierter Fehler- und Retry-Zustand von
    \texttt{RUN-000001}. Originaldokumente:
    \texttt{EVT-000003.json} und \texttt{EVT-000004.json};
    vollständiger Evidenzauszug siehe
    Anhang~\ref{app:traceability_930444_workstation}.}
    \label{fig:result_930444_failed_attempt}
\end{figure}

Der Processing State des ersten Attempts wird damit nicht nachträglich in
einen erfolgreichen Zustand geändert. Der anschließende Retry erhält mit
\texttt{ATT-000002} eine eigene Attempt-Identität, während der vorherige
Fehlerzustand erhalten bleibt.

Neben dem Event State wurde auch die rohe Modellantwort des fehlgeschlagenen
Attempts persistiert. Ein darin enthaltener Interpretationseintrag ist in
Abbildung~\ref{fig:result_930444_raw_failure} auszugsweise dargestellt.

\begin{figure}[htbp]
    \centering
    \begin{minipage}{0.94\textwidth}
        \begin{evidencebox}
{
  "source_evidence_id": "EVD-000001",
  "information_type": "definition",
  "statement_modality": "descriptive",
  "epistemic_class": "explicit",
  "missing_evidence": null,
  "extracton_rationale": "...",
  "uncertainties": [...]
}
        \end{evidencebox}
    \end{minipage}
    \caption{Auszug aus der persistierten Raw Response des
    fehlgeschlagenen Attempts. Originaldokument:
    \texttt{agent\_legacy\_literal\_interpreter\_run\_01.json};
    vollständiger Evidenzauszug siehe
    Anhang~\ref{app:traceability_930444_workstation}.}
    \label{fig:result_930444_raw_failure}
\end{figure}

Der für diesen Verarbeitungsschritt verwendete Parsing-Contract erwartet das
Feld \texttt{extraction\_rationale}. Die persistierte Antwort enthält
dagegen \texttt{extracton\_rationale}. Der zurückgegebene Zustand entspricht
damit nicht dem erwarteten Datenvertrag und wird nicht in den nachgelagerten
Engineering Information Flow übernommen. Der im Event State gespeicherte
Reason Code lautet allgemein
\texttt{shared\_evidence\_interpretation\_failed}; die konkrete
Schemaabweichung bleibt über die unveränderte Raw Response rekonstruierbar.


\paragraph{Source-grounded Evidence und Canonical Engineering Subject}

Im erfolgreichen zweiten Attempt wird die ursprüngliche Source-Aussage zur
\emph{Microscope Workstation} als Evidence Item persistiert. Die zugrunde
liegende Architecture Source enthält die Aussage:

\begin{figure}[htbp]
    \centering
    \begin{minipage}{0.94\textwidth}
        \begin{evidencebox}
The Microscope Workstation provides local session control.
        \end{evidencebox}
    \end{minipage}
    \caption{Source-Aussage zur \emph{Microscope Workstation}.
    Originaldokument:
    \texttt{legacy/demo/safe\_demo/05\_architecture\_notes.md};
    vollständige Testquelle siehe
    Anhang~\ref{app:traceability_930444_inputs}.}
    \label{fig:result_930444_source}
\end{figure}

Die Aussage wird als \texttt{EVD-000002} an
\texttt{SRC-000005}, die zugehörige Source Projection und den konkreten
Source-Segment-Bezug gebunden. Aus dieser Evidence wird anschließend das
Canonical Engineering Subject \texttt{SUBJ-000001} mit dem Label
\emph{Microscope Workstation} gebildet.

Damit referenzieren die nachfolgenden Persona-Schritte denselben
Engineering-Gegenstand. Die einzelnen Perspektiven erzeugen an dieser Stelle
keine voneinander unabhängigen Subject-Identitäten.


\paragraph{Persona Interpretation und Consensus/Variance}

Für \texttt{SUBJ-000001} wurden drei Persona-Interpretationen persistiert.
Bei der Dimension \texttt{information\_type} unterscheiden sich die
Klassifikationen. Der resultierende Consensus-Zustand ist in
Abbildung~\ref{fig:result_930444_consensus} dargestellt.

\begin{figure}[htbp]
    \centering
    \begin{minipage}{0.94\textwidth}
        \begin{evidencebox}
canonical_subject_id: SUBJ-000001

information_type:
  confidence:                low
  consensus_level:           divergent
  review_attention_required: true

Literal Interpreter:
  function

Systems Engineering Interpreter:
  logical_element

Skeptical Ambiguity Interpreter:
  physical_element

uncertainty:
  The exact system boundary and whether this is hardware,
  software, or a combined workstation are not specified.
        \end{evidencebox}
    \end{minipage}
    \caption{Persona-Interpretationen und Consensus-Zustand für
    \texttt{SUBJ-000001}. Originaldokument:
    \texttt{subject\_consensus.json};
    vollständiger Evidenzauszug siehe
    Anhang~\ref{app:traceability_930444_workstation}.}
    \label{fig:result_930444_consensus}
\end{figure}

Technisch enthält dieses Artefakt weiterhin die einzelnen
Persona-Ergebnisse. Für die betrachtete Dimension wird keine Interpretation
automatisch als autorisierter Engineering-Zustand ausgewählt. Stattdessen
wird die Abweichung als \texttt{divergent} sowie
\texttt{review\_attention\_required=true} persistiert und an den Human
Engineering Review übergeben.


\paragraph{Human Engineering Review}

Im Human Engineering Review wird das Canonical Subject als
\texttt{RIT-000016} geführt. Der Reviewer erhält neben dem Subject auch die
zugehörige Source Evidence und die vorbereiteten Interpretations- und
Consensus-Informationen. Der finalisierte Zustand ist auszugsweise in
Abbildung~\ref{fig:result_930444_human_review} dargestellt.

\begin{figure}[htbp]
    \centering
    \begin{minipage}{0.94\textwidth}
        \begin{evidencebox}
RIT-000016 -- Microscope Workstation

Review Outcome:
  accepted_with_modification

Reviewed Content:
  The Microscope Workstation is described as providing
  local session control.

Information Type:
  logical_element

Modality:
  descriptive

Epistemic Status:
  explicit

Human Rationale:
  Logical
        \end{evidencebox}
    \end{minipage}
    \caption{Finalisierter Human-Review-Zustand für
    \emph{Microscope Workstation}. Originaldokument:
    \texttt{reviewed\_report.md};
    vollständiger Review-Auszug siehe
    Anhang~\ref{app:traceability_930444_workstation}.}
    \label{fig:result_930444_human_review}
\end{figure}

Die zugehörige Review Revision speichert zusätzlich die Reviewer-Identität
\texttt{MZ}, die ausgewählten Werte, den Ursprung der Änderung und den
Zeitpunkt der Entscheidung. \texttt{MZ} ist das im Prototyp verwendete
Reviewer-Kürzel des Autors und leitet sich aus \emph{MoritZ} ab.

Der Outcome \texttt{accepted\_with\_modification} kennzeichnet, dass der
vorbereitete Review-Zustand nicht unverändert autorisiert wurde. Als
Information Type wurde \texttt{logical\_element} ausgewählt. Die während des
Live-Testlaufs eingegebene Rationale \texttt{Logical} ist knapp gehalten.
Sie dokumentiert die Entscheidung, den Gegenstand für die weitere
Engineering-Verarbeitung als logisches Element zu behandeln und damit weder
die alternative Klassifikation \texttt{function} noch
\texttt{physical\_element} zu übernehmen.


\paragraph{Approved Engineering Information und Project Fit}

Nach Abschluss des Human Engineering Review wird der freigegebene Zustand als
\texttt{AIN-000011} in Approved Engineering Information überführt. Das
zugehörige Manifest bindet den Inhalt an \texttt{RIT-000016},
\texttt{RVR-000009}, \texttt{SRC-000005}, \texttt{RUN-000001} und
\texttt{ATT-000002}. Neben diesen Referenzen enthält es den autorisierten
fachlichen Zustand mit Titel, Beschreibung, Information Type, Modality und
Epistemic Status.

Für \texttt{SRC-000005} wird anschließend ein eigener Project-Fit-Zustand
persistiert. Dieser erhält im Referenzlauf den Outcome
\texttt{plausible\_in\_scope}. Die zugehörige Rationale hält fest, dass die
beschriebene Microscope Workstation und der Remote Client zum registrierten
Projektkontext der Remote Collaboration in Digital Microscopy passen. Source,
Source Projection, Processing Run und Attempt bleiben dabei weiterhin
referenziert.

Technisch wird durch diesen Schritt kein neues, zusammengeführtes
Approved-Engineering-Information-Artefakt erzeugt. Der source-lokal
autorisierte Zustand bleibt erhalten und wird für die projektweite
Weiterverarbeitung zugelassen.


\paragraph{Human Model Placement und Internal Engineering Model}

Für \texttt{AIN-000011} wird anschließend eine eigene Model Placement
Decision getroffen. Der relevante Zustand ist in
Abbildung~\ref{fig:result_930444_placement} dargestellt.

\begin{figure}[htbp]
    \centering
    \begin{minipage}{0.94\textwidth}
        \begin{evidencebox}
decision_id:       MPD-000016
approved_input_id: AIN-000011
outcome:           accepted
selected_rule_id:  ELEMENT_SYSTEM_LOGICAL
reviewer_identity: MZ
rationale:         Systemlevel
        \end{evidencebox}
    \end{minipage}
    \caption{Human Model Placement Decision für
    \texttt{AIN-000011}. Originaldokument:
    \texttt{MPD-000016.json};
    vollständiger Evidenzauszug siehe
    Anhang~\ref{app:traceability_930444_workstation}.}
    \label{fig:result_930444_placement}
\end{figure}

Die Placement Decision legt die modellbezogene Einordnung der bereits
autorisierten Engineering Information fest. Im resultierenden
\texttt{IEM-000001} wird der Gegenstand als
\texttt{logical\_component} im Bereich \texttt{system.logical} mit der
Framework-Zuordnung \texttt{FW\_SYSTEM\_LOGICAL} materialisiert. Das
Modellelement \texttt{IME-000011} referenziert dabei weiterhin die
Placement Authority \texttt{MPD-000016}.

Die während des Live-Testlaufs eingegebene Rationale
\texttt{Systemlevel} bezeichnet die Entscheidung, den Gegenstand auf der
Systemebene des verwendeten Modellframeworks einzuordnen.


\paragraph{Model Quality Decision und Successor Model}

Im nachfolgenden Model-Quality-Schritt wird für das interne Modellelement eine
modellbezogene Formulierung vorbereitet. Diese wird wiederum als eigene Human
Decision behandelt. Für \texttt{IME-000011} wurde die in
Abbildung~\ref{fig:result_930444_quality} dargestellte Entscheidung
persistiert.

\begin{figure}[htbp]
    \centering
    \begin{minipage}{0.94\textwidth}
        \begin{evidencebox}
decision_id:               MQD-000008
internal_model_element_id: IME-000011
decision:                  approved
reviewer_identity:         MZ

approved_name:
  Microscope Workstation

approved_description:
  Logical component described as providing local
  session control.

rationale:
  Batch-confirmed after the refinement quality contract
  reported preserved meaning, no unsupported information
  and no explicit Human-attention flag.
        \end{evidencebox}
    \end{minipage}
    \caption{Model Quality Decision für
    \texttt{IME-000011}. Originaldokument:
    \texttt{MQD-000008.json};
    vollständiger Evidenzauszug siehe
    Anhang~\ref{app:traceability_930444_workstation}.}
    \label{fig:result_930444_quality}
\end{figure}

Der freigegebene Text wird nicht in den bereits bestehenden
\texttt{IEM-000001}-Snapshot zurückgeschrieben. Stattdessen wird ein
Successor Model \texttt{IEM-000002} erzeugt. Die Elementidentität
\texttt{IME-000011} bleibt erhalten, während die freigegebene Beschreibung
im neuen Modellzustand erscheint.


\paragraph{Ergänzende Authority-Entscheidungen}

Nicht alle im Referenzlauf auftretenden Entscheidungsarten liegen für den
durchgehenden Workstation-Pfad separat vor. Zwei weitere Zustände werden
deshalb ergänzend anhand anderer Gegenstände desselben Referenzlaufs
betrachtet.

Im Review der System Requirements wurde beispielsweise eine vorgeschlagene
Relationship mit der Richtung

\begin{figure}[htbp]
    \centering
    \begin{minipage}{0.94\textwidth}
        \begin{evidencebox}
decision_id:       SRD-000001
source_subject_id: SUBJ-000002
relationship_kind: performs
target_subject_id: SUBJ-000001
outcome:           rejected
rationale:         vice versa
reviewer_identity: MZ
        \end{evidencebox}
    \end{minipage}
    \caption{Beispiel einer verworfenen Relationship Decision.
    Originaldokument: \texttt{SRD-000001.json};
    ergänzende Relationship-Evidence siehe
    Anhang~\ref{app:traceability_930444_additional}.}
    \label{fig:result_930444_relationship}
\end{figure}

verworfen. Eine separat persistierte Proposition mit umgekehrter
Source-/Target-Richtung wurde als \texttt{SRD-000012} akzeptiert. Dabei
handelt es sich um zwei eigenständige Relationship Proposals und nicht um
eine nachträgliche Änderung derselben Decision. Die knappe Rationale
\texttt{vice versa} bezeichnet im Live-Test die Entscheidung, dass die
vorgeschlagene gerichtete Relation fachlich in der Gegenrichtung
weiterverwendet werden soll.

Eine eigene Target-Model-Formulation Decision liegt beispielsweise für den
\emph{Microscope Operator} vor. \texttt{TFD-000001} bindet das interne
Modellelement \texttt{IME-000001} an das Target-Notation-Construct
\texttt{TN\_003} und den Outcome \texttt{materialize\_formally}. Die
persistierte Formulierung lautet:

\begin{figure}[htbp]
    \centering
    \begin{minipage}{0.94\textwidth}
        \begin{evidencebox}
decision_id:                       TFD-000001
authority_subject_id:              IME-000001
selected_relevance_outcome:        materialize_formally
selected_target_notation_construct_id: TN_003

selected_formulation_text:
  part def 'Microscope Operator';

reviewer_identity:
  MZ
        \end{evidencebox}
    \end{minipage}
    \caption{Beispiel einer Target-Model-Formulation Decision.
    Originaldokument: \texttt{TFD-000001.json};
    vollständiger Evidenzauszug siehe
    Anhang~\ref{app:traceability_930444_additional}.}
    \label{fig:result_930444_formulation}
\end{figure}

Der Schritt verändert nicht den zuvor autorisierten fachlichen Gegenstand.
Er persistiert die Entscheidung, wie dieser Gegenstand innerhalb der
zugelassenen Target Notation formal materialisiert werden soll.


\paragraph{Generiertes SysML v2 und Output-Traceability}

Nach Abschluss der autorisierten Modellzustände wird der
\texttt{IEM-000002}-Snapshot deterministisch in SysML v2 serialisiert. Für
den durchgehend betrachteten Gegenstand \emph{Microscope Workstation}
entsteht das in Abbildung~\ref{fig:result_930444_sysml} dargestellte
Modellfragment.

\begin{figure}[htbp]
    \centering
    \begin{minipage}{0.94\textwidth}
        \begin{evidencebox}
part IME_000011 {
    doc /* Engineering name: Microscope Workstation
    Description:
    Logical component described as providing local session control. */
}
        \end{evidencebox}
    \end{minipage}
    \caption{Generiertes SysML-v2-Element für
    \emph{Microscope Workstation}. Originaldokument:
    \texttt{OUT-000001/generated\_model.sysml}, Zeilen 73--77;
    vollständiger Evidenzauszug siehe
    Anhang~\ref{app:traceability_930444_workstation}.}
    \label{fig:result_930444_sysml}
\end{figure}

Parallel zur Generierung wird eine Output-Traceability persistiert. Für
\texttt{IME\_000011} enthält sie den Bezug auf
\texttt{AIN-000011}, die Placement Authority \texttt{MPD-000016}, den
Internal-Model-Zustand \texttt{IEM-000002} sowie die konkrete Position im
generierten Artefakt:

\begin{figure}[htbp]
    \centering
    \begin{minipage}{0.94\textwidth}
        \begin{evidencebox}
approved_input_id:                    AIN-000011
authority_id:                         MPD-000016
authority_type:                       model_placement_decision
generated_symbol_id:                  IME_000011
source_internal_engineering_model_id: IEM-000002
source_internal_model_element_id:      IME-000011

generated_location:
  start_line: 73
  end_line:   77
        \end{evidencebox}
    \end{minipage}
    \caption{Output-Traceability für das generierte Modellelement.
    Originaldokument:
    \texttt{OUT-000001/traceability.json};
    vollständiger Evidenzauszug siehe
    Anhang~\ref{app:traceability_930444_workstation}.}
    \label{fig:result_930444_output_traceability}
\end{figure}

Damit ist die finale Symbolidentität nicht der einzige persistierte Bezug.
Das generierte Element bleibt mit dem autorisierten Engineering Input, seiner
Model Placement Authority und dem tatsächlich serialisierten
Internal-Model-Zustand verbunden.


\paragraph{Technische Validierung und Human Release}

Das generierte SysML-v2-Artefakt wurde anschließend mit dem
SYSIDE Modeler CLI~0.10.3 geprüft. Der persistierte Validation Report weist
für den erzeugten Artifact Set den Status \texttt{valid}, einen
\texttt{passed} Publication Gate und keine Diagnostics oder Findings aus.

\begin{figure}[htbp]
    \centering
    \begin{minipage}{0.94\textwidth}
        \begin{evidencebox}
Project:          930444
Source IEM:       IEM-000002
Validation status: valid
Publication gate:  passed
Findings:          0

External validator:
  SYSIDE Modeler CLI 0.10.3
  execution:   completed
  diagnostics: 0
        \end{evidencebox}
    \end{minipage}
    \caption{Technischer Validierungszustand des finalen
    SysML-v2-Artefakts. Originaldokument:
    \texttt{OUT-000001/validation\_report.md};
    vollständiger Report siehe
    Anhang~\ref{app:traceability_930444_workstation}.}
    \label{fig:result_930444_validation}
\end{figure}

Der Validation State verändert das geprüfte Modell nicht. Nach der
technischen Prüfung wird der exakte Artefaktzustand einem getrennten Final
Model Review zugeführt. Für den Referenzlauf wurden
\texttt{FMR-000001} und die zugehörige Revision
\texttt{FRV-000001} persistiert.

Die anschließende Human Release Decision \texttt{FRD-000001} trägt den
Outcome \texttt{approved\_for\_publication} und die Reviewer-ID
\texttt{MZ}. Gleichzeitig referenziert sie den geprüften
Artifact-Set-Fingerprint und den konkreten Validation-Result-Fingerprint.
Für diese Entscheidung wurde im Live-Test keine zusätzliche Rationale
eingetragen; das entsprechende Feld ist im Originalartefakt
\texttt{null}.

Erst der so gebundene Zustand wird als immutable Published Output
\texttt{OUT-000001} persistiert. Das Output-Paket enthält das generierte
SysML-v2-Modell, die Generation Summary, die Output-Traceability sowie den
maschinen- und menschenlesbaren Validierungszustand.

Der exemplarisch dargestellte Informationspfad lässt sich damit ausgehend
von der Source-Aussage über Evidence, Canonical Subject,
Persona-Interpretationen, Human Engineering Review, Approved Engineering
Information, Model Placement und die internen Modellzustände bis zum
generierten SysML-v2-Symbol verfolgen. Die ergänzenden Beispiele zeigen
zusätzlich die Persistenz einer Relationship Decision und einer expliziten
Target-Model-Formulation Decision.

Die Darstellung bildet bewusst keinen vollständigen Dump des Runtime-
Verzeichnisses ab. Vollständig wiedergegeben werden im Anhang die kurzen
synthetischen Eingangsdokumente und der kompakte Validation Report. Für
umfangreichere Maschinenartefakte werden dort die zur Rekonstruktion des
gezeigten Pfads erforderlichen Originalfelder zusammengeführt. Der Umfang
und die Auswahl des im Anhang dokumentierten Evidenzsatzes sind in
Anhang~\ref{app:traceability_930444_scope} beschrieben.

Die Einordnung, welche Bedeutung die persistierten Zwischenzustände,
Human Decisions, Rationale-Informationen und Authority-Bindungen für die
Kontrollierbarkeit des KI-gestützten Engineering-Prozesses besitzen, erfolgt
davon getrennt in Kapitel~\ref{sec:discussion}.

\subsection{Technische Validierung, Human Release und Ergebnisgrenzen}
\label{sec:results_validation}

Die abschließende technische Validierung des im Referenzlauf erzeugten
SysML-v2-Artefakts wurde mit dem SYSIDE Modeler CLI in Version~0.10.3
durchgeführt. Die Verarbeitung wurde mit Exit Code~0 abgeschlossen; für das
geprüfte Artefakt wurden keine Diagnostics und keine Validation Findings
ausgegeben. Der persistierte Validation State weist entsprechend den Status
\texttt{valid} und einen bestandenen technischen Publication Gate aus.

Dieser technische Zustand wurde nicht unmittelbar als veröffentlichter
Output behandelt. Das validierte Artefakt wurde anschließend einem getrennten
Final Model Review zugeführt. Für den geprüften Modellstand wurden
\texttt{FMR-000001} und die zugehörige Revision \texttt{FRV-000001}
persistiert. Die darauf folgende Human Release Decision
\texttt{FRD-000001} wurde durch den Reviewer \texttt{MZ} mit dem Outcome
\texttt{approved\_for\_publication} abgeschlossen. Erst dieser Zustand wurde
als \texttt{OUT-000001} publiziert.

Der publizierte Output ist dabei an den konkret geprüften Modellzustand
gebunden. Neben dem generierten SysML-v2-Artefakt enthält das Output-Paket die
Generation Summary, die zugehörige Traceability sowie den maschinen- und
menschenlesbaren Validierungszustand. Änderungen des zugrunde liegenden
Modellartefakts würden entsprechend einen neuen Validierungs- und
Freigabezustand erfordern.

Die Ergebnisse der durchgeführten Evaluation beziehen sich auf den
implementierten Proof of Concept und die verwendeten kontrollierten
Testdaten. Die automatisierte Softwareverifikation zeigt, dass die geprüften
deterministischen Contracts und Zustandsübergänge im untersuchten
Implementierungsstand ausführbar waren. Die formative End-to-End-Erprobung
dokumentiert die beobachteten Verarbeitungspfade sowie diejenigen
Architekturgrenzen, die während der verbundenen Ausführung präzisiert werden
mussten. Die Multi-Source-Verifikation zeigt für den untersuchten Testfall,
dass mehrere source-lokal autorisierte Informationspfade in einen gemeinsamen
projektweiten Modellbildungsprozess überführt werden konnten. Der
Referenzlauf~930444 dokumentiert schließlich einen vollständigen
Verarbeitungspfad bis zum technisch validierten und menschlich freigegebenen
SysML-v2-Artefakt.

Aus diesen Ergebnissen wird keine allgemeine Vollständigkeit der
Engineering-Information abgeleitet. Die Verarbeitung und Traceability
beziehen sich auf die im jeweiligen Projekt registrierten und für die
Verarbeitung zugelassenen Informationsquellen. Informationen, die in diesen
Quellen nicht enthalten oder aus ihnen nicht ableitbar sind, werden durch die
Evaluation nicht als vorhanden nachgewiesen.

Ebenso beschreibt der erfolgreiche SYSIDE-Validierungszustand ausschließlich
die technische Verarbeitbarkeit des erzeugten SysML-v2-Artefakts innerhalb
der verwendeten Werkzeugumgebung. Er stellt keine eigenständige fachliche
Bestätigung der im Modell enthaltenen Engineering-Aussagen dar. Diese
fachliche Authority verbleibt bei den im Workflow persistierten Human
Engineering und Model Decisions.

Die durchgeführten Testläufe stellen darüber hinaus keine statistische
Evaluation der Qualität generativer Modelle und keine unabhängige
Usability-Studie dar. Der formative Test wurde durch den Entwickler und Autor
der Arbeit selbst durchgeführt. Beobachtungen zur Bedienbarkeit,
Statusdarstellung und zum manuellen Review-Aufwand beziehen sich daher auf die
prototypische Softwareausprägung und werden nicht als eigenständiger Nachweis
des Architekturkonzepts gewertet.

Damit liegen für die weitere Diskussion drei voneinander getrennte
Ergebnisebenen vor: die technische Verifikation der implementierten
Mechanismen, die beobachtete Ausführung des kontrollierten
Engineering-Information-Flows und die technische Validierung sowie
menschliche Freigabe des resultierenden Modellartefakts. Die Einordnung
dieser Ergebnisse im Hinblick auf die Forschungsfragen, den Stand der Technik
und die Grenzen des entwickelten Architekturkonzepts erfolgt in
Kapitel~\ref{sec:discussion}.